%% Методические указания
%% Лаборатория математики ФН1
\documentclass[12pt, a4paper]{article}
\usepackage[T2A]{fontenc}
\usepackage[utf8]{inputenc}
\usepackage[russian]{babel}
\usepackage[left=25mm, right=20mm, top=20mm, bottom=20mm]{geometry}
\usepackage{amsmath, amssymb}
\usepackage{enumitem}
\pagestyle{plain}
\setlength{\parindent}{1.25cm}
\begin{document}
\begin{center}
  {\large\bfseries Методические указания}

  \vspace{0.3em}
  {\large Исследование числовых рядов на сходимость}

  \vspace{0.5em}
  {\small Математический анализ · 1 курс\\ Лаборатория математики ФН1}
\end{center}
\vspace{1em}

\section*{Порядок действий}

Проверку сходимости ряда $\sum a_n$ удобно вести по одной и той же схеме.
Переходите к следующему пункту, только если предыдущий не дал ответа.

\begin{enumerate}[nosep]
  \item \textbf{Необходимый признак.} Если $a_n \not\to 0$, ряд расходится, и на этом всё.
  \item \textbf{Знакопостоянный ряд?} Если да~--- пункты 3--5. Если знакочередующийся~--- пункт 6.
  \item \textbf{Сравнение.} Подберите эталон: $\sum 1/n^p$ сходится при $p>1$,
        геометрический $\sum q^n$~--- при $|q|<1$.
  \item \textbf{Признак Даламбера} при наличии факториалов и степеней:
        $\lim |a_{n+1}/a_n| = L$. При $L<1$ сходится, при $L>1$ расходится, при $L=1$ неинформативен.
  \item \textbf{Признак Коши} при наличии $n$-х степеней: $\lim \sqrt[n]{a_n} = L$.
  \item \textbf{Признак Лейбница} для знакочередующегося ряда: модули монотонно убывают и стремятся к нулю.
        Отдельно проверьте абсолютную сходимость ряда $\sum |a_n|$.
\end{enumerate}

\section*{Типичные ошибки}

\begin{itemize}[nosep]
  \item Из $a_n \to 0$ делают вывод о сходимости. Это неверно: гармонический ряд расходится.
  \item Признак Даламбера применяют при $L=1$ и делают вывод. При $L=1$ признак не работает~---
        нужен другой.
  \item Путают сходимость и абсолютную сходимость. Ряд $\sum (-1)^n/n$ сходится, но не абсолютно.
  \item При сравнении берут эталон «на глаз», не проверяя неравенство для всех $n$ начиная с некоторого.
\end{itemize}

\section*{Разобранный пример}

Исследуем $\displaystyle\sum_{n=1}^{\infty} \frac{n!}{n^n}$.

Необходимый признак ответа не даёт. Ряд знакопостоянный, есть факториал~---
применяем Даламбера:
\[
\frac{a_{n+1}}{a_n} = \frac{(n+1)!}{(n+1)^{n+1}}\cdot\frac{n^n}{n!}
= \frac{n^n}{(n+1)^n} = \left(\frac{n}{n+1}\right)^{n} \to \frac{1}{e} < 1.
\]
Ряд сходится.

\end{document}
